\chapter{How to read this atlas}

This atlas collects every architecture and workflow diagram of the project in one place. It begins with the models
\method{} is built from, redrawn from their original papers: the YOLO family from its first version to YOLO26, and the
state-space models that lead to Mamba. It then draws \method{} itself, the way it is trained and evaluated, and the
system it is meant to grow into. Every number in the diagrams of \method{} was read from the trained model or from
its configuration, not estimated.

\section{One visual language}

All diagrams share the conventions of the Transformer paper~\cite{vaswani2017attention}: rounded blocks with a
light fill and a thin outline, grey containers marked $N\times$ for blocks that repeat, small circles for
element-wise operations, and thin arrows for data. Diagrams of whole networks follow the layout of recent
multispectral detection papers such as CFMW~\cite{li2025cfmw}: dashed regions group the parts of the network,
real images and feature maps sit where the data flows, and tensor shapes are written beside the layers that
produce them.

Colour always means the same thing.

\begin{tabularx}{\textwidth}{@{}P{3.4cm}L@{}}
\hd{colour} & \hd{meaning} \\
blue & the visible camera, its backbone and its feature maps \\
orange & the thermal camera and its backbone \\
green & fusion by weighted sum, the neck, and outputs that are kept \\
purple & selective scans and state: the Mamba block, CMFM and the temporal memory \\
gold & detection heads, attention blocks and evaluation \\
pink & linear projections and embeddings \\
brown & data and the steps that prepare it \\
grey & convolutions and plain bookkeeping such as concatenation and upsampling \\
\end{tabularx}

\section{Arrows, shapes and symbols}

A solid arrow carries data. A dashed arrow carries something that exists only in training, is optional, or is
planned. The circles mark element-wise operations: $+$ adds, $\times$ multiplies, $\Vert$ concatenates along the
channel axis and $\div$ splits it. A shape written as $128\times80^2$ means 128 channels on an 80 by 80 grid, for a
640-pixel input; the stride of a level is the input size divided by its grid size.

\section{What is redrawn and what is ours}

Chapters~2 and~3 redraw published architectures in the common style above. They follow the original papers in
structure and naming, and they use this project's images only to make the data flow concrete. Chapters~4 to~6 are
the project's own: \method{}, its training and evaluation, and its roadmap. Appendix~A lists the source of every
figure.
